\section{Henkin semantics}
\label{sec:second-order-logic-05}

\begin{definition}[Henkin structure]
A Henkin structure for the relational second-order language is a structure
\[
  \mathcal H=(M,(\mathcal R_n)_{n\geq1},\ldots),
\]
where \(M\ne\varnothing\), each
\(\mathcal R_n\) is a nonempty subset of \(\Pow(M^n)\), and the remaining
components interpret the nonlogical symbols.  An \(n\)-ary relation variable
ranges over \(\mathcal R_n\), not over all of \(\Pow(M^n)\).
\end{definition}

A systematic treatment of Henkin semantics and of the translation between
second-order logic and many-sorted first-order logic is
\cite{Manzano1996}.

When the calculus contains a comprehension schema, the semantic class
must be restricted accordingly: every relation definable by an eligible
formula with parameters from the Henkin domains must belong to the relevant
\(\mathcal R_n\).  This closure is a hypothesis on the semantic class.  It is
not automatic from \(\mathcal R_n\subseteq\Pow(M^n)\).

\begin{definition}[Henkin semantic consequence]
Write
\[
  \Gamma\entails_{\mathrm H}\varphi
\]
when every Henkin structure under every sorted assignment satisfying
\(\Gamma\) satisfies \(\varphi\).  A full second-order structure is the
special case
\(\mathcal R_n=\Pow(M^n)\) for every \(n\).
\end{definition}

\paragraph{Cited theorem.}
For the typed formula system and general structures fixed in
\cite[pp.~81--91]{Henkin1950}, a formula is a theorem of that source's system
iff it is valid in that source's general structures.  The hypotheses
include the source's typed axiom and rule surface and its corresponding closure
conditions; the conclusion is completeness for that paired syntax and
semantics, not for full second-order validity.  This is a cited theorem.
For Church's 1940 typed system, general models, an exact first-order
encoding and a proved completeness theorem are given in
\cref{sec:p02-w02-p02-040-church-general}.  It
does not define derivability for this chapter, whose relational language has
not imported that rule surface or selected a comprehension schema.

\paragraph{Worked many-sorted translation.}
The semantic reduction underlying Henkin completeness can nevertheless be
stated exactly without choosing a proof calculus.  Use a many-sorted
first-order language with an object sort \(D\), a relation sort \(S_n\) for
each \(n\geq1\), and an \((n+1)\)-ary predicate symbol \(E_n\) whose argument
sorts are \((S_n,D,\ldots,D)\).  All original constants, function symbols, and
relation symbols keep their object-sort interpretations.

To distinguish a second-order formula from its many-sorted image, define the
metalinguistic translation \(\varphi\mapsto\varphi^\sharp\) recursively.  It
leaves object terms, equality, and nonlogical relation atoms unchanged,
commutes with every propositional connective, and has the clauses
\[
\begin{aligned}
  \bigl(X^n(t_1,\ldots,t_n)\bigr)^\sharp
    &\defeq E_n(X,t_1,\ldots,t_n),\\
  (\forall x\,\psi)^\sharp
    &\defeq\forall x{:}D\,\psi^\sharp,
  & (\exists x\,\psi)^\sharp
    &\defeq\exists x{:}D\,\psi^\sharp,\\
  (\forall X^n\,\psi)^\sharp
    &\defeq\forall X{:}S_n\,\psi^\sharp,
  & (\exists X^n\,\psi)^\sharp
    &\defeq\exists X{:}S_n\,\psi^\sharp.
\end{aligned}
\]
This translation is not mere renaming: it replaces relation application by an
incidence predicate and makes the two quantifier ranges into sorts.  Its
standing is a recursively eliminable metalinguistic operation, and it is used
in the satisfaction equivalence below.

From a Henkin structure
\(\mathcal H=(M,(\mathcal R_n)_{n\geq1},\ldots)\), form the many-sorted
structure \(\mathcal H^\sharp\) by taking
\[
  D^{\mathcal H^\sharp}=M,
  \qquad
  S_n^{\mathcal H^\sharp}=\mathcal R_n,
  \qquad
  E_n^{\mathcal H^\sharp}(R,\bar a)
  \Longleftrightarrow \bar a\in R.
\]
This structure is needed because \(\varphi^\sharp\) contains sorted variables
and the incidence symbols \(E_n\), which the original second-order structure
does not interpret as first-order symbols.  Its standing is a many-sorted
first-order structure.  It is nonredundant with \(\mathcal H\): relations in
\(\mathcal H\) are values of second-order variables, whereas the same
relations are first-order objects of the sorts \(S_n\) in
\(\mathcal H^\sharp\).  The structure is used immediately in the following
equivalence.
For every second-order formula \(\varphi\) and every matching sorted assignment
\(s\),
\[
  \mathcal H,s\entails\varphi
  \quad\Longleftrightarrow\quad
  \mathcal H^\sharp,s\entails\varphi^\sharp.
\]
Indeed, induction on \(\varphi\) handles an atom \(X^n(\bar t)\) by the
definition of \(E_n\); first-order atoms because the object interpretations are
unchanged; connectives by the induction hypothesis; object quantifiers because
both sides range over \(M\); and relation quantifiers because both sides range
over \(\mathcal R_n\).  These cases exhaust the formula grammar.

If a selected calculus contains a comprehension instance for an eligible
formula \(\theta(x_1,\ldots,x_n,\bar z,\bar Y)\), where \(X^n\) is not free
in \(\theta\), its semantic counterpart is the many-sorted sentence
\[
  \forall\bar z\,\forall\bar Y\,
  \exists X{:}S_n\,\forall\bar x{:}D\,
  \bigl(E_n(X,\bar x)\leftrightarrow
        \theta^\sharp(\bar x,\bar z,\bar Y)\bigr),
\]
where \(\bar x=(x_1,\ldots,x_n)\), and every displayed parameter is quantified
in its declared sort.  Thus the corresponding closure condition says exactly
that the \(n\)-ary relation defined by each eligible instance belongs to
\(\mathcal R_n\); it is not automatic for an arbitrary family of relation
domains.

\paragraph{Internal proof boundary.}
The translation proves a semantic equivalence, not a local completeness
theorem.  Before this chapter may assert Henkin completeness, compactness, or a
downward L\"owenheim--Skolem result for its own calculus, it must still list
every axiom, rule, structural convention, equality and extensionality
principle, and eligible comprehension instance; define finite derivations;
match those instances with semantic closure; and carry the chosen many-sorted
first-order completeness construction back to exactly that rule surface.  No
calculus name or derivability--validity biconditional is introduced here.

Henkin completeness and full categoricity concern different consequence
relations.  Combining them into the assertion ``second-order logic is complete
and categorical'' suppresses the semantic parameter and is incorrect.

